% Section 4
\section{Districts and the rung ladder}\label{sec:districts}

``District'' means any partition of the units into groups: ridings, counties
within states, schools within boards. Nothing below uses the nesting being
geographic. Let $D\in\{1,\dots,D^\star\}$ with $P(D=d)>0$, $V=(X,N,D)$, and let
$A_d$ be a compact convex set with $P(X\in A_d\mid D=d)=1$; the default is the
district's own composition hull. Write
$\mu^{(d)}_{rs}=\E[NX_rX_s\one\{D=d\}]$ and $\mu^{(d)}_r=\E[NX_r\one\{D=d\}]$, so
that $\mu_r=\sum_d\mu^{(d)}_r$. For a vector $\zeta\in\R^R$, $\zeta\one^\top$ is the
matrix whose row $r$ is constant and equal to $\zeta_r$; since $\one^\top x=1$,
$(\zeta\one^\top)x=\zeta$, so this is an additive shift written in homogeneous form.

The district model replaces (S2) by $\E[\beta\mid V]=M_{H^{(D)}}(X)$ and (S3) by
validity on each $A_d$, with one set of coefficients
$H^{(d)}=(H^{(d)}_1,\dots,H^{(d)}_C)$ per district. Write $Q^{(d)}$ and $K^{(d)}$
for the symmetric and skew parts and $k^{(d)}$ for the free coordinates of
$K^{(d)}$. The rungs are restrictions on $(H^{(1)},\dots,H^{(D^\star)})$, listed in
decreasing order of plausibility:

\begin{center}\small
\begin{tabularx}{\linewidth}{@{}L{0.14\linewidth}Y@{}}
\toprule
Rung & Restriction \\
\midrule
M2 & none: each district has its own linear contextual effects \\
M1($\delta$) & there is $\bar k$ with $\|k^{(d)}-\bar k-c_d\|_\infty\le\delta$ for all $d$, where $c_d$ are the coordinates of $\skw(\zeta_{\cdot,d}\one^\top)$ and $\zeta_{c,d}=\diag(Q^{(d)}_c-Q^{(1)}_c)$ \\
M1 & $H^{(d)}_c=H_c+\zeta_{c,d}\one^\top$ with $\sum_c\zeta_{c,d}=0$ and $\zeta_{c,1}=0$: shared contextual effects, district baselines \\
L4($r$) & M1, and row $r$ has no contextual effect: $H_{c,rs}=H_{c,rr}$ for all $s,c$ \\
L5 & $H^{(d)}_c=\alpha_{c,d}\one^\top$: constancy within each district (Goodman by district) \\
L6 & L5 with $\alpha_{c,d}$ common to all districts (Goodman) \\
M0 & one $H$ for all units on $\conv\bigcup_dA_d$: the pooled model of Section~\ref{sec:rxc}, not nested with the others \\
\bottomrule
\end{tabularx}
\end{center}

The parameter $\delta$ is in probability units. Entries of $K$ are half-gaps
between two hypothetical rates, so $\delta=0.05$ lets the gap between
``$r$-members in pure-$s$ units'' and ``$s$-members in pure-$r$ units'' differ
across districts by up to 10 points.

\subsection{Grouping alone adds nothing}

\begin{proposition}[No model, no gain]\label{prop:D0}\status{proved}
Under (S1) alone, the sharp identified set for $B^*_{rc}$ is the Duncan--Davis
interval $\DD^*_{rc}$ of Proposition~\ref{prop:nonempty}, whatever the law of $D$
or of any other covariate. The joint sharp set for the matrix $B^*$ is convex and
compact but in general not the product of the cell intervals; for a linear
target the sharp interval is
$\big[-\E[N\,h_S(-\tilde a(X);X,T)],\ \E[N\,h_S(\tilde a(X);X,T)]\big]$.
\end{proposition}
\begin{reading}
District dummies cannot narrow the model-free bound, which is already computed
unit by unit. Every gain from districts comes from an assumption that links
them. The accounting identity links the cells within each unit, so the per-cell
extremes cannot in general be attained together, and a combined target can be
strictly narrower than the sum of its cells' intervals.
\end{reading}

\subsection{What each rung identifies}

\begin{lemma}[Identification under the district rungs]\label{lem:D1}\status{proved}
Assume A5 within each district.
(a) Under every rung, $Q^{(d)}$ is identified from $\E[T\mid X,D=d]=X^\top Q^{(d)}X$,
$\sum_cH^{(d)}_c=\one\one^\top$, and each $K^{(d)}\in\Kspace$ is unidentified.
(b) Under M1, $\zeta_{c,d}=\diag(Q^{(d)}_c-Q^{(1)}_c)$ is identified; the
off-diagonal entries satisfy the overidentifying restrictions
$(Q^{(d)}_c-Q^{(1)}_c)_{rs}=\tfrac12(\zeta_{c,d,r}+\zeta_{c,d,s})$ for $r\neq s$;
$K^{(d)}=K+\skw(\zeta_{\cdot,d}\one^\top)$ with $K=K^{(1)}$ common, and the
unidentified dimension equals that of the pooled model.
(c) Under L5, $\alpha_{c,d}=\diag Q^{(d)}_c$, $K^{(d)}=\skw(\alpha_{\cdot,d}\one^\top)$
is identified, and the structure $Q^{(d)}_c=\sym(\alpha_{c,d}\one^\top)$ is testable.
(d) Under M1, A5 within every district can be replaced by a stacked rank
condition: the design made of the pooled quadratic monomials and the district
shifts $\one\{D=d\}X_r$ ($d\ge2$) has full column rank in the population. A
sufficient condition is A5 in one district and $R$ affinely independent
compositions in every other district; $R$ distinct but affinely dependent
compositions are not enough.
\end{lemma}
\begin{reading}
District baselines are identified from each district's own regression. Under
shared effects the curvature of the regression must be the same in every
district: that is the testable part, with $(D^\star-1)(C-1)R(R-1)/2$
restrictions. The shared contextual split $K$ is not testable. Given the common
$Q$, a district's regression is linear in its shift, so a district needs only a
few compositions to fix its baseline, for example two distinct units in
$2\times2$; but with three groups they must not lie on a line. This is
identification, not precision: with about ten units per district the shifts are
identified but noisy.
\end{reading}

\begin{lemma}[The estimand under districts]\label{lem:D2}\status{proved}
Under (S1) and the district version of (S2), with $P(X\in A_D)=1$,
\[
B^*_{rc}=\frac1{\mu_r}\sum_d\Big[\sum_sQ^{(d)}_{c,rs}\mu^{(d)}_{rs}+\sum_{s\ne r}K^{(d)}_{c,rs}\mu^{(d)}_{rs}\Big].
\]
Under M1 its sensitivity to the common $K_{c,rs}$ is $\mu_{rs}/\mu_r$, as in the
pooled model.
\end{lemma}
\begin{reading}
The electorate-wide rate is a weighted sum of district rates. Under shared
effects the unknown enters exactly as in the pooled model, so the district
baselines change the known part of the target, not how much the unknown matters.
\end{reading}

\subsection{Sharp and outer sets for each rung}

\begin{theorem}[Sharpness for affine mean models]\label{thm:affine}\status{proved}
Let a parameter $\theta$ range over a closed set, and let $\mathcal V_A=\{v=(x,n,d):x\in A_d\}$.
Suppose (H1) $v\mapsto m_\theta(v)\in\R^{R\times C}$ is Borel on $\mathcal V_A$;
(H2) $\sup_{v\in\mathcal V_A}\|m_\theta(v)\|<\infty$; (H3) $\theta\mapsto m_\theta(v)$
is affine, and continuous uniformly in $v$ on bounded sets of $\theta$; (H4) the
unit polytopes are not affected by $\theta$. Define the model by (S1),
$\E[\beta\mid V]=m_\theta(V)$ on $\{V\in\mathcal V_A\}$, and $m_\theta(v)\in\Bset$ on
$\mathcal V_A$. Then $\theta$ is in the sharp set iff the validity condition holds
and $\langle u,m_\theta(v)\rangle\le\psi(u;v)$ for $P_V$-almost every $v$ and all $u$.
\end{theorem}
\begin{reading}
The sharpness argument of Theorem~\ref{thm:sharp} carries over to any model in
which the mean is affine in the unknowns, including every district rung.
\end{reading}

\begin{theorem}[The rungs and the ladder]\label{thm:ladder}\status{proved}
Assume A1, A5 within districts and $P(X\in A_D)=1$.
(a) For each rung, the sharp set is the rung's restriction intersected with each
district's own sharp condition (Theorem~\ref{thm:sharp} within district $d$).
(b) The outer set is
$\{\theta\in\text{rung}:(Q^{(d)}_c+K^{(d)}_c(\theta))v\ge0\ \forall d,c,\ v\in\vertx(A_d)\}$.
It is a polyhedron, and every target bound is an LP. For M1 the variables are the
common $k$ and $\Theta^{M1}_O=\bigcap_d\Theta^{(d)}_O$, where $\Theta^{(d)}_O$
imposes district $d$'s rows with its identified shift.
(c) Write $\mathcal B^\rho$ for the set of values of a linear target under rung
$\rho$. For $0\le\delta\le\delta'$,
\[
\mathcal B^{L6}\subseteq\mathcal B^{L5}\subseteq\mathcal B^{L4(r)}\subseteq\mathcal B^{M1}\subseteq\mathcal B^{M1(0)}\subseteq\mathcal B^{M1(\delta)}\subseteq\mathcal B^{M1(\delta')}\subseteq\mathcal B^{M2}\subseteq\DD^*,
\]
for both the sharp and the outer versions, whenever the smaller rung's model is
non-empty.
\end{theorem}
\begin{reading}
Each stronger assumption gives a bound inside the weaker one's, so the ladder can
be reported as a table and readers can see what each assumption buys. The pooled
model M0 is not on this chain. M1(0) keeps the common skew part but drops the
testable symmetric restrictions of Lemma~\ref{lem:D1}(b); it is a sensitivity
benchmark that separates a failure of M1's testable implications from
uncertainty about the unidentified effects, not a substantive rung.
\end{reading}
\begin{sketch}
Each rung is an affine mean model, so Theorem~\ref{thm:affine} applies district
by district. The outer version replaces the sharp condition by non-negativity at
the vertices of each district hull. The nesting holds because each step restricts
the same polyhedral system further; L5 is M1 with constant rows.
\end{sketch}

\subsection{Closed forms and properties}

\begin{proposition}[Shared effects with two rows]\label{prop:D3}\status{proved}
Let $R=2$ and $\vertx(A_d)=\{(t,1-t):t\in\{l_d,u_d\}\}$ with $0<l_d\le u_d<1$, and
let $\lambda_c,\rho_c$ be as in Proposition~\ref{prop:tworows} with $Q=Q^{(1)}$.
Under M1,
\[
\ell_c=\max_d\max_{t\in\{l_d,u_d\}}\Big[\lambda_c(t)-\frac{\zeta_{c,d,1}}{1-t}\Big],\qquad
h_c=\min_d\min_{t\in\{l_d,u_d\}}\Big[\rho_c(t)+\frac{\zeta_{c,d,2}}{t}\Big],
\]
and $\Theta^{M1}_O=\{\kappa:\ell\le\kappa\le h,\ \sum_c\kappa_c=0\}$, projected by
Lemma~\ref{lem:box}. If a district has a pure vertex ($u_d=1$ or $l_d=0$), the
corresponding row at that vertex does not involve $\kappa_c$: it is a testable
restriction and is dropped from the maximum (or minimum), while the other row at
that vertex enters as the formula gives.
\end{proposition}
\begin{reading}
A district raises the lower bound when its baseline for group 1 is low and it has
units that are almost all group 1: there, group 1's rate is nearly observed, and
the shared contextual effect carries that information to every district. One
district with lopsided units and an extreme baseline can tighten the bound for
the whole electorate. The upper bound behaves symmetrically through group 2.
\end{reading}
\begin{sketch}
Under M1, district $d$'s mean for row 1 at vertex $t$ is
$Q_{c,11}t+(Q_{c,12}+\kappa_c)(1-t)+\zeta_{c,d,1}$. Requiring it to be non-negative
and dividing by $1-t>0$ gives a lower limit on $\kappa_c$; row 2 gives an upper
limit after dividing by $t$. The shared $\kappa_c$ must satisfy every district's
limits at every vertex, and the sum-to-zero constraint across columns then gives
a box cut by a hyperplane.
\end{sketch}

\begin{proposition}[When district baselines matter]\label{prop:D4}\status{proved}
(a) If every shift is zero, M1's outer set equals the pooled outer set on
$\conv\bigcup_dA_d$. (b) Under M1 with $P(X\in A_D)=1$, write $Z_d=[\zeta_{c,d,r}]_{rc}$.
Then $\E[\beta\mid X,N]=M_H(X)+\E[Z_D\mid X,N]$, and the pooled model M0 holds
(for some coefficients) iff $\E[Z_D\mid X,N]$ is almost surely a function of $X$
alone that is affine on $\supp P_X$. Independence of $D$ and $(X,N)$ suffices, and
so does conditional independence of $D$ and $N$ given $X$ together with
$\E[Z_D\mid X]$ affine in $X$; independence of $D$ and $X$ alone does not. When M0
holds in this way, its validity condition (S3) is automatic on the population hull
$\conv\supp P_X$ (Lemma~\ref{lem:outer}(c)); on any other pooled region it remains a
separate condition.
\end{proposition}
\begin{reading}
The pooled model is safe when district baselines, given composition and unit
size, vary at most linearly with composition. When strongholds also have
different baselines (a popular local candidate in the strongest districts), or
when district baselines correlate with unit size, the pooled fit mistakes
baseline differences for contextual effects and can miss the truth: an ecological
omitted-variable bias. If two districts have the same distribution of
compositions but one has larger units and a non-zero shift, $\E[Z_D\mid X,N]$
varies with $N$ and the size-weighted pooled target is biased, although the
$X$-only pooled model still holds. Section~\ref{sec:tests} tests the stronger,
design-conditional condition that all shifts vanish.
\end{reading}

\begin{proposition}[District-specific effects]\label{prop:D5}\status{proved}
Under M2 the outer bound is the Minkowski sum of the district bounds,
$\mathcal B^{M2}_O=\sum_d[\varphi^{(d)}_L,\varphi^{(d)}_U]$, where $\varphi^{(d)}$ is
district $d$'s term in Lemma~\ref{lem:D2}. The same holds for the sharp version.
\end{proposition}
\begin{reading}
District-specific lines are safe but rarely much tighter than Duncan--Davis,
because widths add across districts.
\end{reading}

\begin{proposition}[Bounded heterogeneity]\label{prop:D6}\status{proved}
(a) $\delta\mapsto\mathcal B^{M1(\delta)}_O$ is nested. (b) For
$\delta\ge\bar\delta:=\max_d\max_{k\in\Theta^{(d)}_O}\|k-c_d\|_\infty$, which is
finite, it equals $\mathcal B^{M2}_O$. (c) Its endpoints are continuous and
piecewise linear in $\delta$ wherever it is non-empty. (d) At $\delta=0$ it equals
the M1 outer bound computed with the same $Q^{(d)}$ and shifts.
\end{proposition}
\begin{reading}
The dial $\delta$ traces a continuous path between district-specific and shared
effects, so a reader can see how much the shared-effects assumption is doing.
\end{reading}

\begin{proposition}[The most restrictive rungs]\label{prop:D7}\status{proved}
(a) L5 point-identifies the target:
$B^*_{rc}=\sum_d(\mu^{(d)}_r/\mu_r)\,\alpha_{c,d,r}$ with $\alpha_{c,d}=\diag Q^{(d)}_c$,
Goodman's regression within each district. If L5 holds, this point lies in every
larger rung's bound. (b) L4($r$) pins $K_{c,rs}=Q_{c,rr}-Q_{c,rs}$ (reference
district) for all $s\ne r$ and all $c$, which is $(C-1)(R-1)$ of the free
coordinates; for $R=2$ the target is point-identified.
\end{proposition}
\begin{reading}
The constancy rungs give single numbers. Reported beside the bounds, they show
what ``no contextual effects'' implies. If row $r$ does have a contextual effect,
L4($r$) attributes it to the other rows, and the point ($R=2$) or the bound is
misspecified and can exclude the truth, with no warning except an empty set.
With non-voters as row $r$ this is a real risk, because mobilisation often
depends on local campaign intensity. If some district's contextual split differs
from the shared one by more than $\delta$, M1($\delta$) fails in the same way as
M1, but less sharply: the larger $\delta$, the more protection and the less
narrowing.
\end{reading}

\subsection{Worked examples}\label{sec:ladder-examples}

The examples are noise-free, so they concern identification only.

\paragraph{Two groups, two districts.} For in-party loyalty with truth 0.819:
Duncan--Davis gives $[0.453,0.879]$; the pooled model is empty, because the
districts have different baselines; M2 gives $[0.563,0.870]$; and M1 gives
$[0.618,0.866]$. The L5 point is 0.858, which is misspecified, since the true
contextual slopes are 0.10.

\paragraph{Three groups, six districts.} The districts have different
composition ranges and baselines, with a common $H$:

\begin{center}\small
\setlength{\tabcolsep}{3.5pt}
\begin{tabular}{@{}lcccccc@{}}
\toprule
Target & Truth & DD & M2 & M1(.05) & M1 & L5 \\
\midrule
Loyalty P$\to$P & 0.836 & [0.167, 0.953] & [0.391, 0.916] & [0.397, 0.914] & [0.430, 0.880] & 0.870 \\
Defection P$\to$O & 0.095 & [0.000, 0.610] & [0.024, 0.404] & [0.024, 0.390] & [0.042, 0.339] & 0.061 \\
Demobilisation P$\to$N & 0.068 & [0.000, 0.551] & [0.024, 0.350] & [0.024, 0.323] & [0.044, 0.295] & 0.069 \\
Mobilisation N$\to$P & 0.135 & [0.001, 0.685] & [0.045, 0.470] & [0.051, 0.444] & [0.089, 0.413] & 0.122 \\
\bottomrule
\end{tabular}
\end{center}

The ladder shrinks monotonically and contains the truth at every rung. Moving
from M2 to M1 narrows the loyalty bound from width 0.525 to 0.450, and
M1($\delta$) at $\delta=0.10,0.05,0.02$ traces a continuous path between them. The
baseline shifts here are small (standard deviation 0.04); the gain grows with more
heterogeneous baselines and more lopsided units. The L5 point for loyalty is
0.870 against a truth of 0.836: inside the M1 bound, but wrong. The pooled model
is misspecified here, because the baselines vary with composition; it happens to
cover the truth, but Proposition~\ref{prop:D4} gives no guarantee that it will.
